\documentclass[11pt,a4paper]{article}
\usepackage[margin=2.4cm]{geometry}
\usepackage{amsmath,amsthm,mathtools}
\usepackage{fontspec}
\usepackage{unicode-math}
\setmainfont{STIX2Text}[Path=fonts/, Extension=.otf,
  UprightFont=*-Regular, ItalicFont=*-Italic,
  BoldFont=*-Bold, BoldItalicFont=*-BoldItalic]
\setmathfont{STIX2Math}[Path=fonts/, Extension=.otf]
\usepackage[protrusion=true,expansion=false]{microtype}
\usepackage{booktabs,enumitem}
\usepackage{placeins}
\usepackage{tikz}
\usetikzlibrary{arrows.meta,positioning,calc,fit,backgrounds}
\definecolor{tierspend}{RGB}{200,60,60}
\definecolor{tierfull}{RGB}{210,140,40}
\definecolor{tierin}{RGB}{60,130,180}
\definecolor{tierrecv}{RGB}{70,150,90}
\usepackage[psdextra,hidelinks]{hyperref}

\emergencystretch=3em
\hbadness=10000
\hfuzz=2pt

\newtheorem{theorem}{Theorem}[section]
\newtheorem{lemma}[theorem]{Lemma}
\newtheorem{proposition}[theorem]{Proposition}
\newtheorem{corollary}[theorem]{Corollary}
\theoremstyle{definition}
\newtheorem{definition}[theorem]{Definition}
\newtheorem{construction}[theorem]{Construction}
\newtheorem{assumption}[theorem]{Assumption}
\newtheorem{remark}[theorem]{Remark}
\newtheorem{example}[theorem]{Example}

% kind-coded theorem blocks, palette shared with the figures and the website:
% definitions/constructions blue (tierin), assumptions amber (tierfull),
% proved statements a neutral bar; remarks and examples run unframed
\usepackage{mdframed}
\providecolor{tierin}{RGB}{60,130,180}
\providecolor{tierfull}{RGB}{210,140,40}
\mdfdefinestyle{kindbar}{topline=false,bottomline=false,rightline=false,
  linewidth=1.5pt,innerleftmargin=9pt,innerrightmargin=4pt,
  innertopmargin=5pt,innerbottommargin=6pt,leftmargin=0pt,rightmargin=0pt,
  skipabove=10pt,skipbelow=10pt}
\surroundwithmdframed[style=kindbar,linecolor=tierin!70,backgroundcolor=tierin!4]{definition}
\surroundwithmdframed[style=kindbar,linecolor=tierin!70,backgroundcolor=tierin!4]{construction}
\surroundwithmdframed[style=kindbar,linecolor=tierfull!80,backgroundcolor=tierfull!5]{assumption}
\surroundwithmdframed[style=kindbar,linecolor=black!30]{theorem}
\surroundwithmdframed[style=kindbar,linecolor=black!30]{lemma}
\surroundwithmdframed[style=kindbar,linecolor=black!30]{proposition}
\surroundwithmdframed[style=kindbar,linecolor=black!30]{corollary}

\usepackage{fancyhdr}
\pagestyle{fancy}
\fancyhf{}
\fancyhead[L]{\footnotesize\scshape The Zcash Arboretum}
\fancyhead[R]{\footnotesize\itshape\nouppercase{\leftmark}}
\fancyfoot[C]{\footnotesize\thepage}
\renewcommand{\headrulewidth}{0.3pt}
\renewcommand{\sectionmark}[1]{\markboth{\thesection\ \ #1}{}}

\newcommand{\F}{\mathbb{F}}
\newcommand{\Z}{\mathbb{Z}}
\newcommand{\N}{\mathbb{N}}
\newcommand{\G}{\mathbb{G}}
\newcommand{\E}{\mathbb{E}}
\newcommand{\PP}{\mathbb{P}}
\renewcommand{\Pr}{\mathrm{Pr}}
\newcommand{\Fp}{\F_p}
\newcommand{\Fq}{\F_q}
\newcommand{\ip}[2]{\langle #1,\, #2 \rangle}
\newcommand{\Comm}{\mathsf{Commit}}
\newcommand{\Open}{\mathsf{Open}}
\newcommand{\Setup}{\mathsf{Setup}}
\newcommand{\Verify}{\mathsf{Verify}}
\newcommand{\negl}{\mathsf{negl}}
\newcommand{\poly}{\mathsf{poly}}
\newcommand{\PPT}{\mathsf{PPT}}
\newcommand{\Adv}{\mathcal{A}}
\newcommand{\Prov}{\mathcal{P}}
\newcommand{\Ver}{\mathcal{V}}
\newcommand{\Sim}{\mathcal{S}}
\newcommand{\Ext}{\mathcal{E}}
\newcommand{\Rel}{\mathcal{R}}
\newcommand{\Lang}{\mathcal{L}}
\newcommand{\nfsym}{\mathsf{nf}}
\newcommand{\cmsym}{\mathsf{cm}}
\newcommand{\cmx}{\mathsf{cmx}}
\newcommand{\cvsym}{\mathsf{cv}}
\newcommand{\rk}{\mathsf{rk}}
\newcommand{\Pallas}{\ensuremath{\mathsf{Pallas}}}
\newcommand{\Vesta}{\ensuremath{\mathsf{Vesta}}}
\newcommand{\Ep}{\ensuremath{\mathcal{E}}}
\newcommand{\NoteCommit}{\ensuremath{\mathsf{NoteCommit}}}
\newcommand{\Sinsemilla}{\ensuremath{\mathsf{Sinsemilla}}}
\newcommand{\ShortCommit}{\ensuremath{\mathsf{ShortCommit}}}
\newcommand{\CommitIvk}{\ensuremath{\mathsf{Commit}^{\mathsf{ivk}}}}
\newcommand{\Extract}{\ensuremath{\mathsf{Extract}}}
\newcommand{\Acc}{\ensuremath{\mathsf{Acc}}}
\newcommand{\GroupHash}{\ensuremath{\mathsf{GroupHash}}}
\newcommand{\ask}{\ensuremath{\mathsf{ask}}}
\newcommand{\aksym}{\ensuremath{\mathsf{ak}}}
\newcommand{\nksym}{\ensuremath{\mathsf{nk}}}
\newcommand{\ivk}{\ensuremath{\mathsf{ivk}}}
\newcommand{\fvk}{\ensuremath{\mathsf{fvk}}}
\newcommand{\ovk}{\ensuremath{\mathsf{ovk}}}
\newcommand{\dk}{\ensuremath{\mathsf{dk}}}
\newcommand{\pkd}{\ensuremath{\mathsf{pk_d}}}
\newcommand{\gd}{\ensuremath{\mathsf{g_d}}}
\newcommand{\rcv}{\ensuremath{\mathsf{rcv}}}
\newcommand{\rcm}{\ensuremath{\mathsf{rcm}}}
\newcommand{\rivk}{\ensuremath{\mathsf{r_{ivk}}}}
\newcommand{\rsk}{\ensuremath{\mathsf{rsk}}}

\renewenvironment{abstract}{%
  \small\begin{center}{\bfseries\abstractname}\end{center}\medskip\noindent\ignorespaces
}{\par\medskip}

\title{\textbf{\Huge FROST Guide}\\[6pt]\large Threshold custody for randomised RedPallas keys}
\author{m@rek.onl}
\date{}
\begin{document}
\maketitle
\thispagestyle{empty}
\begin{abstract}
This volume of \emph{The Zcash Arboretum} constructs threshold spend
authorisation: several holders of key shares produce one ordinary
RedPallas signature without reconstructing the spending key during
signing. It assumes interpolation from the \emph{Math Guide}, Schnorr
and RedPallas from the \emph{Crypto Guide}, Action authorisation from the
\emph{Ironwood Guide}, and separated transaction-building roles from the
\emph{Wallet Guide}. The construction includes share provisioning,
transcript encodings, nonce state, randomisation and the final signing
workflow. Algebraic correctness, imported unforgeability results and
unproved composition requirements are distinguished explicitly.
\end{abstract}
\clearpage
\tableofcontents
\clearpage

\section{Threshold custody and its contract}\label{sec:fr-context}

\emph{Math Guide}, \S\emph{Groups and scalar multiplication} and
\S\emph{Prime fields and modular arithmetic}, supplies the group and
scalar language used here. \emph{Crypto Guide},
\S\emph{Security contracts and reductions}, defines signature
unforgeability and the adversarial experiments used to state it.

A spending key normally belongs to one signer. Threshold custody replaces
that signer by \(n\) participants with a threshold \(t\): any authorised
set of at least \(t\) participants can cooperate, while fewer than \(t\)
should not be able to forge a signature. The chain still sees one public
key and one signature. It does not learn the threshold or the selected
participants from an additional on-chain list of signatures.

This volume selects the signing mechanism of
\href{https://www.rfc-editor.org/rfc/rfc9591.html}{RFC 9591}, Sections 4--5,
and its Pallas integration in \texttt{frost} at
\texttt{2016e44b} and \texttt{reddsa} at \texttt{3792daa9}.
The exact randomizer construction is the seed-and-commitments API in
\texttt{frost-rerandomized/\allowbreak src/\allowbreak lib.rs}; the group
and hash functions are in
\texttt{reddsa/\allowbreak src/\allowbreak frost/\allowbreak redpallas.rs}.
\href{https://zips.z.cash/zip-0312}{ZIP 312} supplies the Zcash
compatibility requirements, but its direct-randomizer helper is not
silently substituted for this selected implementation.

The endpoint is threshold authorisation of an Orchard-style Action.
No Crosslink, voting or future recovery protocol is needed to understand
that endpoint.

\subsection{The selected group, encodings, and hashes}

A \emph{ciphersuite} is a fixed choice of group, encodings, and hash
functions. \emph{Crypto Guide}, \S\emph{Hash functions and key derivation},
constructs the BLAKE2b operation used in this suite. We fix the
functions before using them in the protocol equations.

The group is Pallas, with order \(q=p_{\Vesta}\) and base
\(G=G^{\mathsf{SpendAuth}}\), as constructed in \emph{Crypto Guide},
\S\emph{RedPallas signatures}. Canonical compressed points occupy
32 bytes. Scalar encoding is a 32-byte little-endian integer strictly
less than \(q\); decoding rejects an out-of-range integer rather than
reducing it. Participant identifiers additionally reject zero.

The group decoder rejects malformed points and the identity at FROST
input positions that require non-identity elements. Algebraic identities
below still hold when a group sum is the identity, but the concrete
ciphersuite's serialisation and validation errors govern whether that
rare session can be completed. It must fail cleanly, not reinterpret an
invalid encoding.

For a 64-byte string \(b\), write
\(\operatorname{Wide}(b)=\operatorname{LEInteger}(b)\bmod q\).
The scalar-valued hashes apply this map to BLAKE2b-512's output.
The complete personalised functions are:

\begin{center}
\begin{tabular}{@{}lll@{}}
\toprule
Function & Output & BLAKE2b personalisation\\
\midrule
\(H_1\) & scalar & \texttt{FROST\_RedPallasR}\\
\(H_2\) & scalar & \texttt{Zcash\_RedPallasH}\\
\(H_3\) & scalar & \texttt{FROST\_RedPallasN}\\
\(H_4\) & 64 bytes & \texttt{FROST\_RedPallasM}\\
\(H_5\) & 64 bytes & \texttt{FROST\_RedPallasC}\\
\(H_R\) & scalar & \texttt{FROST\_RedPallasA}\\
\(H_D\) & scalar & \texttt{FROST\_RedPallasD}\\
\bottomrule
\end{tabular}
\end{center}

The hash \(H_D\) supports participants jointly generating shares of
a key without giving any participant the complete secret. It is an
additional library function, not an RFC signing hash.
Identifiers may simply be assigned as distinct nonzero scalars;
no identifier hash is needed for that choice. We write
\(P^\star\) for the canonical encoding of a point \(P\), and
\(\operatorname{enc}(x)\) for a scalar or identifier encoding.
The selected key-generation
API takes 16-bit participant counts and requires \(2\leq t\leq n\),
although the mathematical sharing construction also makes sense at
\(t=1\).


\subsection{Public and secret objects}

Work in an additive prime-order group \(\G=\langle G\rangle\) of order
\(q\). Scalars belong to \(\F_q\). The Pallas instantiation will set
\(q=p_{\Vesta}\). The \emph{Crypto Guide}, \S\emph{Knowledge, simulation,
and Schnorr}, constructs the verification equation
\[
 [z]G=R+[c]Y,\qquad c=H_2(R^\star\|Y^\star\|m).
\]
Here \(Y=[x]G\) is the group public key, \(m\) is the approved byte-string
message, \(R\) is the nonce commitment and \(z\) the response scalar.
FROST distributes the calculation of \(z\); it does not change this
outside verification rule.

Each participant has a distinct nonzero identifier \(i\in\F_q\),
a secret signing share \(x_i\), and a public verification share
\(Y_i=[x_i]G\). A public key package fixes \(Y\), all identifiers and
their \(Y_i\), and the threshold. A signer does not accept a coordinator's
replacement verification-share map as a new definition of its key.

The coordinator collects nonce commitments, distributes one signing
package and aggregates response shares. Its ability to schedule or
interrupt the protocol does not give it spending authority. Conversely,
unforgeability does not oblige a malicious coordinator or selected signer
to finish a session.

\begin{definition}[Threshold authorisation requirement]
\label{def:fr-contract}
Fix an honestly provisioned key and an adversary controlling fewer than
\(t\) participants. Honest participants approve the complete application
message before answering a signing request. A successful forgery is an
accepting signature outside the signing interactions permitted by the
security game's freshness rule: the rule specifying which earlier
signing interactions disqualify a proposed forgery. The game must state
when participants may be corrupted, whether sessions may overlap,
who chooses each key randomizer, and the exact disqualifying interactions.
\end{definition}

These qualifications matter. A new signature on a previously signed
message is excluded by an ordinary new-message game but included by a
strong-unforgeability game. A theorem for corruption choices fixed before
key generation does not automatically cover secrets exposed adaptively
during signing. We identify the imported guarantees in
\S\ref{sec:fr-security}.

\section{Shares and verifiable key generation}\label{sec:fr-shamir}

The interpolation theorem is proved in \emph{Math Guide},
\S\emph{Polynomials and evaluation}. We use it to construct a secret
sharing and then to explain precisely what the public share checks add.
That same section proves the hidden-constant-coefficient lemma used
for the privacy statement below.

\subsection{A polynomial with the key at zero}

Choose \(1\leq t\leq n<q\), a secret \(x\in\F_q\), and independent
uniform coefficients \(a_1,\ldots,a_{t-1}\in\F_q\). Define
\[
 f(X)=x+\sum_{j=1}^{t-1}a_jX^j,\qquad x_i=f(i).
\]
The degree is at most \(t-1\); a highest coefficient may be zero.
Identifiers cannot be zero because \(f(0)\) is the secret itself.

For a signing set \(S\) of distinct identifiers with \(|S|\geq t\),
define
\[
 \lambda_i^S=\prod_{\substack{j\in S\\j\ne i}}
                  \frac{j}{j-i}.
\]
All denominators are invertible because identifiers are distinct.
Evaluation of the interpolation formula at zero gives
\[
 \sum_{i\in S}\lambda_i^Sx_i=x,
 \qquad
 \sum_{i\in S}\lambda_i^S=1.
\]
The second equality is interpolation of the constant polynomial \(1\).
It will be essential when all shares receive the same randomizer.

This construction is called \emph{Shamir secret sharing}. Any fixed set
of \(r<t\) private shares is jointly uniform in \(\F_q^r\), independently
of the fixed secret \(x\): apply the lemma ``A hidden constant coefficient''
in \emph{Math Guide}, \S\emph{Polynomials and evaluation}, with its
constant \(s=x\) and evaluation points equal to the participant identifiers.
Its requirements are exactly the nonzero distinct identifiers and independent
uniform nonconstant coefficients enforced here.

Publishing \(Y=[x]G\) changes the information-theoretic statement:
an unbounded observer can recover \(x\) from \(Y\). Threshold key secrecy
with public verification data is computational. The private-share statement
must not be quoted as perfect privacy of the entire public-key protocol.

\subsection{Checking a dealer's shares}

A trusted dealer can distribute the shares over confidential,
authenticated channels and publish coefficient commitments
\[
 C_j=[a_j]G\quad(0\leq j<t),\qquad a_0=x.
\]
Each recipient checks
\[
 [x_i]G=\sum_{j=0}^{t-1}[i^j]C_j.
\]
The right side equals \([f(i)]G\). Since scalar multiplication by a
generator is injective over \(\F_q\), the equality verifies that the
received scalar is the evaluation described by those public commitments.

All participants must receive the same commitment list. Otherwise a
dealer could give individually consistent but mutually incompatible
views. An authenticated private channel identifies the sender to one
recipient; it does not by itself establish agreement between recipients.

The dealer knows \(x\). Splitting an existing key is therefore compatible
with a seed-derived account, but its threshold guarantee requires the
original full key and dealer randomness to be erased or otherwise
accounted for. Dividing a key does not erase previously made backups.

\subsection{Generating shares without a single dealer}

The selected library also implements distributed key generation (DKG) in
\texttt{frost-core/\allowbreak src/\allowbreak keys/\allowbreak dkg.rs}.
It runs a dealer contribution from
every participant and adds the results. We state its essential messages,
including the checks needed to obtain one common key.

Each participant \(\ell\) samples a polynomial
\(f_\ell(X)=\sum_{j=0}^{t-1}a_{\ell j}X^j\). In the selected
implementation the constant coefficient is sampled nonzero, as a
signing key; the other coefficients are independent uniform scalars.
It publishes
\(C_{\ell j}=[a_{\ell j}]G\). It proves knowledge of its constant
coefficient using a Schnorr proof, with a fresh nonzero nonce
\(u_\ell\):
\[
 U_\ell=[u_\ell]G,\quad
 e_\ell=H_D(\operatorname{enc}(\ell)\|C_{\ell0}^\star\|U_\ell^\star),
 \quad
 v_\ell=u_\ell+e_\ell a_{\ell0}.
\]
The verifier recomputes \(e_\ell\) and checks
\([v_\ell]G=U_\ell+[e_\ell]C_{\ell0}\).
The selected implementation hashes exactly these three fields.
Its hash function separates the DKG role; an identifier for this one
key-generation run, called a ceremony identifier, and
membership agreement remain responsibilities of the enclosing transport
protocol, not extra inputs we invent for this library call.

After agreeing on all first-round packages, participant \(\ell\) sends
\(f_\ell(i)\) privately to recipient \(i\). Each recipient verifies
\[
 [f_\ell(i)]G=\sum_{j=0}^{t-1}[i^j]C_{\ell j}
\]
for every contributor. It rejects missing, duplicate, unauthenticated
or inconsistent packages. On success it forms
\[
 x_i=\sum_{\ell=1}^n f_\ell(i),\quad
 Y_i=\sum_{\ell=1}^n\sum_{j=0}^{t-1}[i^j]C_{\ell j},\quad
 Y=\sum_{\ell=1}^n C_{\ell0}.
\]
The sum \(F(X)=\sum_\ell f_\ell(X)\) has degree at most \(t-1\),
\(x_i=F(i)\), and \(Y=[F(0)]G\). Thus the shares have exactly the
interpolation relation signing needs.

The private contributions and polynomial coefficients are erased after
their required checks and aggregation. The enduring private object is
the participant's resulting key share. The enduring public object is
the agreed key package.

This is an aborting DKG, not a proof that malicious participants cannot
block key generation. A proof of knowledge of each contribution also
does not by itself prove unbiased output when an adversary aborts
unfavourable runs and allows only favourable runs to complete.
The signing theorem below does not silently acquire a DKG-composition
guarantee merely because the final interpolation identities are correct.

\subsection{Normalising an Orchard authorisation key}

Orchard's long-term spend-authorisation key uses the even-\(y\)
representative; \emph{Ironwood Guide}, \S\emph{Keys and viewing capabilities},
defines the convention. If the generated \(Y\) has odd \(y\), negate
every secret share and corresponding public share, the group key, and
any retained polynomial commitments:
\[
 x_i'=-x_i,\qquad Y_i'=-Y_i,\qquad Y'=-Y.
\]
Then
\(\sum_i\lambda_i^Sx_i'=-x\), so the new shares interpolate to
the normalised group key. Negating only the group key would leave an
inconsistent key package.

The selected RedPallas \texttt{post\_generate} and \texttt{post\_dkg}
hooks apply this transformation. The later per-Action randomised key
has no even-\(y\) requirement. Applying the normalisation again to that
key would change the transaction's authorisation equation.

\section{Two-round signing}\label{sec:fr-sign}

We first construct signing for one fixed group key \(Y\). The later
randomisation step changes the key inputs but preserves the equations.

\subsection{Nonce generation and first-round commitments}

For each signing session, participant \(i\) creates two nonce scalars
\(d_i,e_i\), called the hiding and binding nonces. The RFC's hedged
nonce generator combines fresh randomness with the secret share:
\[
 \operatorname{Nonce}(x_i;b)=H_3(b\|\operatorname{enc}(x_i)),
 \qquad b\gets\{0,1\}^{256}.
\]
Invoke it independently twice, using two fresh 32-byte strings, to obtain
\(d_i,e_i\). The participant retains the scalars and sends only
\[
 D_i=[d_i]G,\qquad E_i=[e_i]G
\]
to the coordinator. Hashing the secret share into nonce generation does
not make repeated random inputs safe. Identical inputs repeat the nonce.

The participant associates the pair with durable single-use state.
Publishing a commitment is not permission to reuse its secret in several
responses. Preprocessing may prepare pairs before the message is known;
each pair is still consumed by at most one signing operation.

\subsection{The canonical signing package}

The coordinator selects \(S\), with \(t\leq|S|\leq n\), and constructs
\[
 B=((i,D_i,E_i))_{i\in S}
\]
in ascending numerical identifier order. The group encoding is
\[
 \operatorname{Enc}(B)=
 \mathop{\|}_{i\in S}
 \bigl(\operatorname{enc}(i)\|D_i^\star\|E_i^\star\bigr).
\]
Within one ciphersuite all three encodings have fixed lengths; this
makes the concatenation unambiguous. Ascending scalar order is not
lexicographic order of little-endian byte strings: that order compares
the first differing byte, not the integers represented by all bytes.

Each signer receives \(m,B\), validates every point and identifier,
checks uniqueness and ordering, checks that the selected participants
belong to the stored key package, and verifies that its own identifier
and exact commitments occur once. It must also approve \(m\) according
to the application; a syntactically valid byte string is not necessarily
an authorised payment.

Define the message and commitment-list prehashes and binding factors:
\begin{align*}
 M&=H_4(m),& C&=H_5(\operatorname{Enc}(B)),\\
 \rho_i&=H_1(Y^\star\|M\|C\|\operatorname{enc}(i)).
\end{align*}
These are RFC 9591's ordered inputs. The public key, message, complete
commitment list and identifier all participate. The symbol \(\rho_i\)
here denotes a FROST binding factor, not a note's uniqueness seed.

Each participant computes
\[
 R_i=D_i+[\rho_i]E_i,\qquad
 R=\sum_{i\in S}R_i,\qquad
 c=H_2(R^\star\|Y^\star\|m).
\]
No secret nonce is needed to recompute these public values.

\subsection{Responses and aggregation}

Participant \(i\) calculates the Lagrange coefficient for this exact
\(S\), and returns
\[
 z_i=d_i+\rho_i e_i+c\lambda_i^Sx_i.
\]
The coordinator checks scalar encodings, the response-to-identifier
mapping and completeness of the selected response set. It computes
\[
 z=\sum_{i\in S}z_i,\qquad \sigma=R^\star\|\operatorname{enc}(z).
\]
The public verification equation is the ordinary Schnorr equation.
A coordinator may first verify the aggregate and verify individual
shares only after aggregate failure. To check a share, it tests
\[
 [z_i]G=R_i+[c\lambda_i^S]Y_i.
\]
These checks use the stored \(Y_i\), not a public share supplied alongside
an allegedly bad response. Authenticated messages make a failed check
attributable to its sender; the equation alone does not prove who sent
an arbitrary byte string.

\begin{theorem}[Signing correctness]\label{thm:fr-correct}
If all selected shares belong to the same key polynomial, all
participants use the same valid signing package, and all computed
values satisfy the ciphersuite's encoding and nonidentity checks,
the aggregate verifies under \(Y\). Freshness of the nonce pairs is
additionally required for security, not for this algebraic identity.
\end{theorem}
\begin{proof}
Summing the response equations and using interpolation gives
\[
 z=\sum_{i\in S}(d_i+\rho_i e_i)+
      c\sum_{i\in S}\lambda_i^Sx_i
   =r+cx,
 \qquad r=\sum_{i\in S}(d_i+\rho_i e_i).
\]
The group commitment is \(R=[r]G\). Multiplying the scalar equality
by \(G\) yields \([z]G=R+[c]Y\).
\end{proof}

A signer may leave before answering. The coordinator cannot simply drop
it and sum the others: the participant set is inside the binding factors
and interpolation weights, and the dropped commitments are inside \(R\).
A retry selects fresh nonce pairs and constructs a fresh signing package.

\subsection{A scalar-field walkthrough}

For an arithmetic illustration only, take \(q=101\),
\(f(X)=42+17X+29X^2\), and \(S=\{1,3,5\}\).
The shares and coefficients are
\[
 (x_1,x_3,x_5)=(88,51,44),\qquad
 (\lambda_1,\lambda_3,\lambda_5)=(65,24,13).
\]
Both \(\sum_i\lambda_i=1\) and \(\sum_i\lambda_ix_i=42\) hold modulo
101. Choose illustrative nonce pairs
\[
 (d_1,d_3,d_5)=(11,13,17),\qquad
 (e_1,e_3,e_5)=(19,23,29).
\]
For illustrative hash outputs
\((\rho_1,\rho_3,\rho_5)=(2,4,6)\) and \(c=7\), the effective nonces
are \((49,4,90)\), whose sum is \(42\). The response shares are
\((93,88,54)\); their sum is \(33\). Finally,
\[
 33=42+7\cdot42\pmod{101}.
\]
The retained numerical check recomputes every entry. These invented
hash outputs and small scalars test the algebra; they are not a
ciphersuite test vector or an instance with cryptographic security.

\section{Security and nonce state}\label{sec:fr-security}

Correctness says honest responses add to a valid signature. Security
must explain why a malicious coordinator interacting with several
sessions cannot obtain an unapproved one. The equations alone do not
establish that conclusion.
The random-oracle model, which treats each distinct hash input as
receiving one independently random answer consistently reused on
repetition, is constructed in \emph{Crypto Guide},
\S\emph{Security contracts and reductions}.

\subsection{What the binding factor prevents}

A response is linear in the secret quantities:
\[
 z_i=d_i+\rho_i e_i+c\lambda_i^Sx_i.
\]
If the coordinator could select the coefficients of this equation
independently after collecting nonce commitments, it could try to combine
responses from separate sessions into an equation it wanted. The
binding factor makes the effective nonce depend on the entire signing
package, including the eventual message and public key. Changing another
participant's commitment therefore changes the honest participant's
coefficient too.

The domain-separated hashes are essential here. Calling two different
functions the same random oracle, or omitting the participant list from
one side's calculation, changes the protocol analysed by the security
proof. The fact that the resulting aggregate still resembles a Schnorr
signature is not a repair.

\subsection{The imported signing guarantee}

The signing-layer result we rely on is the analysis of FROST1 in
Bellare, Crites, Komlo, Maller, Tessaro and Zhu,
\emph{Better than Advertised Security for Non-interactive Threshold
Signatures}, CRYPTO 2022; the FROST analysis is retained in
\href{https://eprint.iacr.org/2022/833}{ePrint 2022/833}, Theorem 5.4.
Its setting uses random oracles, trusted-dealer key distribution and
static corruption, meaning the corrupted set is fixed before key
generation, of fewer than \(t\) participants. Its strong
threshold-forgery notion also accounts for which participants
contributed to a session; an ordinary single-signer theorem against
new-message forgery after chosen signing queries is not that result.

The hardness assumption is \emph{one-more discrete logarithm}. In its
game, an adversary obtains independent random challenge group elements
and may ask a discrete-log oracle for selected logarithms. It must not
be able to output the logarithms of more challenge elements than the
number of oracle queries it used, with non-negligible probability.
This assumption permits oracle access that the ordinary discrete-log
game does not. We do not equate the two assumptions.

The imported result covers the matching signing construction under
its stated key-generation and corruption hypotheses. It does not prove
that the implemented DKG composes with signing, that arbitrary adaptive
corruptions are safe, or that backup software preserves the nonce
requirements. Those are separate obligations. The purpose of importing
a theorem is to state a precise guarantee, not to replace the
unforgeability question with an adjective.

\subsection{Why a nonce pair has only one response}

If one nonce pair is answered under three packages, the attacker obtains
equations
\[
 z_j=d+\rho_j e+\beta_jx_i,\qquad
 \beta_j=c_j\lambda_i^{S_j},\qquad j=1,2,3.
\]
Subtract the first equation from the other two. This leaves two linear
equations in \(e,x_i\), solved by the elimination procedure in
\emph{Math Guide}, \S\emph{Linear algebra and uniform masks}. Whenever
\[
 (\rho_2-\rho_1)(\beta_3-\beta_1)
 -(\rho_3-\rho_1)(\beta_2-\beta_1)\ne0,
\]
the coefficient matrix is invertible and reveals \(x_i\), then \(d\).
Fewer observations can suffice in special cases or with extra leaked
information. A failed aggregate does not make an emitted response
harmless: the attacker still possesses its linear equation.

Nonce generation from a deterministic function of just the share and
message is therefore unsuitable. Signing the same message in another
coordinator-chosen package would repeat the secret pair while changing
its public coefficients. Fresh randomness and single-use storage solve
different parts of this problem; neither substitutes for the other.

\begin{construction}[Application nonce state]\label{con:fr-noncestate}
Associate each generated pair with a unique local record.
\begin{enumerate}
\item Store the private pair as available before advertising its public
  commitments.
\item Bind that record to one complete accepted signing request.
\item Irreversibly mark it consumed before releasing any response
  derived from it. A crash must not roll that transition back.
\item Erase the private pair when its response is computed or the
  session is abandoned. A restart obtains new commitments.
\end{enumerate}
\end{construction}

This state machine is an application requirement, not a persistence
service furnished by a scalar-signing function. The selected Rust API
takes references to nonce objects; the surrounding application must not
call it again with a consumed pair. Restoring a backup, cloning a device
or racing two signing requests must not duplicate available nonce state.
Ordinary durable key backups should not restore an old pool of
preprocessed nonces.

A malicious participant can withhold a response. Share verification can
identify inconsistent responses from authenticated senders; it cannot
force participation or prove that a sender intended to disrupt the
session. Availability policy and retry policy remain outside the
unforgeability theorem.

\section{Randomised RedPallas signing}\label{sec:fr-rerand}

The \emph{Ironwood Guide}, \S\emph{Balance, digests, and authorisation},
publishes a per-Action key
\[
 \rk=Y+[\alpha]G
\]
rather than the long-term \(Y=\aksym\). The Action proof establishes the
relation to the authorised spend key. The outside signature must verify
under \(\rk\), not under \(Y\).

\subsection{Shifting all shares}

For a common scalar \(\alpha\), define
\[
 \bar x_i=x_i+\alpha,\qquad
 \bar Y_i=Y_i+[\alpha]G,\qquad
 \bar Y=Y+[\alpha]G.
\]
All signers use \(\bar Y\) in the binding factors and signature
challenge and \(\bar x_i\) in their response equation. The coordinator
uses \(\bar Y_i\) for share verification and \(\bar Y\) for final
verification.

\begin{proposition}[Randomisation commutes with sharing]
\label{prop:fr-rerand}
The shifted shares interpolate to \(x+\alpha\) for every permitted
signing set \(S\). Their FROST aggregate is an ordinary signature under
\(\bar Y\).
\end{proposition}
\begin{proof}
The two interpolation identities give
\[
 \sum_{i\in S}\lambda_i^S(x_i+\alpha)
 =x+\alpha\sum_{i\in S}\lambda_i^S=x+\alpha.
\]
Equivalently, the shifted polynomial is \(f(X)+\alpha\), whose degree
remains at most \(t-1\). Theorem~\ref{thm:fr-correct} then applies to the
shifted inputs.
\end{proof}

Adding the same scalar to every share does not add \(|S|\alpha\) to
the group key, because shares are interpolated rather than simply
summed. This is why no separate sharing ceremony for the randomizer is
needed.

\subsection{The selected randomizer construction}

Let \(B\) be the exact first-round commitment list for the Action.
The coordinator draws a 32-byte seed \(s\) and computes
\[
 \alpha=H_R(s\|\operatorname{Enc}(B)).
\]
It sends \(s,B\) to each participant. Each participant regenerates
\(\alpha\) locally and constructs the shifted key package.

This is the mechanism of
\texttt{Randomizer::new\_from\_commitments} and
\texttt{regenerate\_from\_seed\_\allowbreak and\_commitments} at the selected
\texttt{frost} commit. The message is \emph{not} inside this randomizer
hash. It is inside \(H_4\), every binding factor and the signature
challenge. That distinction permits the builder to calculate
\(\rk\) before the transaction digest containing \(\rk\) is known.

Fresh commitments contribute input variation even if the coordinator's
seed generator malfunctions. This is hedging, not a theorem that every
maliciously selected seed produces an unbiased randomizer. A coordinator
can examine candidate seeds and can learn \(\alpha\). Its privacy role
must be accounted for explicitly.

\subsection{Signature compatibility}

Every list entry is therefore \(32+32+32=96\) bytes before transport
framing. The final signature is \(32+32=64\) bytes. Neither number
depends on whether the selected signing set has two participants or
twenty; only the intermediate transcript grows with its size.

The essential compatibility condition is \(H_2\). Its personalisation,
input order, scalar conversion and point/scalar encodings are exactly
RedPallas's. A verifier does not care whether the valid response was
computed from one secret or from interpolated shares. Other FROST
ciphersuites are not interchangeable merely because they also output
64-byte signatures.

\subsection{Privacy and the remaining theorem boundary}

For an independent uniform \(\alpha\), the map
\(\alpha\mapsto Y+[\alpha]G\) is a bijection from \(\F_q\) to \(\G\).
Thus the randomised public key alone is uniform, independently of
\(Y\). With wide reduction of a random-oracle output, the distribution
is statistically close to uniform, rather than exactly uniform.

A person learning \(\alpha\) can compute \(Y=\rk-[\alpha]G\).
The coordinator, signers and transaction prover therefore are not
covered by a claim that their own signing sessions are unlinkable to
the long-term key. The seed must travel over a confidential,
authenticated channel: \(s\) and public \(B\) reveal \(\alpha\).
Transaction timing, signer identity and network traffic may reveal
additional links.

Gouv\^ea and Komlo's
\href{https://eprint.iacr.org/2024/436}{\emph{Re-Randomized FROST}}
provides the algebraic construction and a claimed unforgeability
analysis. Its stated freshness condition concerns new messages; it
does not establish the stronger fresh-message--signature-pair property
requested by the Zcash specification. The selected seed-and-commitments
API also needs its precise correspondence to an analysed randomizer
model. We do not present the preprint's claimed concrete bound as a
proved bound for this full implementation and DKG composition.

The identity in Proposition~\ref{prop:fr-rerand}, consensus verification
compatibility and the public-key masking calculation are proved here.
A full threshold strong-unforgeability and joint privacy theorem for
the complete selected workflow remains a separate security obligation.

\section{Transaction construction with separated signers}
\label{sec:fr-deploy-pczt}

A signer receiving only a 32-byte digest cannot determine what payment
it authorises. Threshold cryptography does not fix a missing transaction
review. The complete signer contract is the one in
\emph{Wallet Guide}, \S\emph{Constructing and authorising a transaction};
this section inserts the threshold rounds at the necessary points.

\subsection{The order that avoids a circular digest}

For each real spend assigned to a threshold key:
\begin{enumerate}
\item Select the participants and obtain their unused first-round
  commitments. At this point the final transaction message need not
  exist.
\item Freeze \(B\), generate \(s\), derive \(\alpha\), and construct
  \(\rk=Y+[\alpha]G\).
\item Insert that \(\rk\) into the Action. Complete the transaction
  effects, including outputs, public balances, fee and all other fields
  covered by the applicable signature digest.
\item Compute the transaction's signature digest \(m\). Give each signer
  the transaction information needed to verify its authorisation policy,
  along with \(m,B,s\) and the identification of the Action it is
  authorising.
\item Each signer checks that its regenerated \(\rk\) equals the Action's
  key, recomputes the required digest, approves the transaction, and
  produces one response using its consumed nonce pair.
\item Aggregate, verify under the Action's \(\rk\), and insert the
  resulting spend-authorisation signature. Complete the proof and
  remaining authorising data in the order their own interfaces permit.
\end{enumerate}

The proof relation uses the same \(\alpha\) that defines the signature
key. The prover must possess the required proof-authorisation data;
threshold spend-key shares alone are not the entire Orchard witness.
Likewise, withholding the spend key from the prover does not make a
malicious builder's proposed outputs safe to approve.

The final signature digest may exclude some authorising data by design.
The signer follows the exact digest definition for that transaction
version; it must not hash an arbitrary serialisation or replace the
digest by the transaction identifier. A partially created transaction
container transports roles and fields; it is not itself proof that the
signer checked them.

\subsection{Applying the shared payment example}

The \emph{Wallet Guide}'s running transaction spends an 80,000-zatoshi
note, pays 50,000, returns 20,000 in change and pays a 10,000 fee.
Its two padded Actions may have net values \(+30{,}000\) and
\(-20{,}000\). FROST changes none of that arithmetic.

The holders of the real note's threshold spend authority jointly sign
its Action. A dummy spend still follows its own permitted authorisation
construction; it does not represent another threshold-owned 80,000
zatoshi note. Before approving, the real signers check recipient,
amount, change ownership, fee and the relevant Action/digest bindings.
The aggregate signature does not reveal which subset approved, but the
participants and coordinator know their own session.

The binding signature is a different role. It supplies the authorisation
component of the balance argument, whose opening-knowledge premises,
proof relations and integer bounds are stated in \emph{Ironwood Guide},
\S\emph{Balance, digests, and authorisation}. Threshold spend
authorisation does not replace that proof of balance or distribute its
key automatically.

\subsection{What a completed endpoint guarantees}

A complete honest run outputs the same signature format the existing
Action verifier accepts, with no reconstruction of the spend key during
signing. The off-chain system must still provide an agreed key package,
participant authentication, private share distribution, confidential
randomizer transport, application-level transaction approval and durable
nonce consumption.

Those requirements are not optional implementation polish: each closes
a specific step in the construction. The threshold-signing endpoint is
complete only when the surrounding custody system implements them.
Nothing in this construction turns a proposed threshold application
for some other protocol into an implemented one.

\end{document}
